% 01_introduction.tex — 第0节 引言（15-20%）
\section{引言}
\label{sec:intro}

随着"双碳"目标推进，风电、光伏等新能源在电力系统中的渗透率持续攀升。国家能源局数据显示，2024年全国风电、光伏装机容量突破12亿千瓦，新能源发电量占比超过XX\%。高比例新能源接入使电力系统呈现"双高"特征——高电力电子化（同步发电机占比降低）和高运行不确定性（出力随机波动），对暂态稳定安全评估提出了新挑战\cite{vittal2022re_stability}。

传统的确定性暂态稳定评估方法针对单一或少量预想工况进行时域仿真，难以覆盖高维运行方式空间。近年来，基于数据驱动的代理模型方法为高效安全评估提供了新途径。然而，现有方法面临三方面不足：

\textbf{（1）单一代理模型难以捕获多约束耦合特征}。暂态稳定受功角、频率、电压多约束共同作用\cite{hatziargyriou2020severity}。现有GP代理模型大多针对单一指标（如功角稳定裕度）建模\cite{wang2023gp_transient}，无法同时预测多约束严重度，难以揭示新能源渗透率变化引起的约束主导模式转换规律\cite{wang2024voltage_re}。Chen等\cite{chen2024mogp_voltage}提出了多输出GP用于电压稳定预测，但未考虑功角和频率约束的耦合。

\textbf{（2）安全边界辨识缺乏高效采样策略}。直接通过网格搜索或蒙特卡洛采样辨识安全域边界计算代价巨大。贝叶斯优化（BO）在超参数优化\cite{snoek2012practical}和实验设计领域已证明采样效率优势，近年来开始应用于电力系统场景选择\cite{bo2025der_scenarios, li2024bo_scenario}和新能源稳定性分析\cite{zhang2023bo_renewable}。但现有BO应用多聚焦于寻找最严重场景（worst-case），而非系统性地辨识完整的安全域边界。

\textbf{（3）边界辨识与采样策略缺乏闭环反馈}。Shahidinejad等\cite{shahidinejad2024gp_bo_transient}将GP与BO结合用于暂态稳定边界探索，但采用一次性采样策略，边界质量受限于初始采样覆盖度。Liu等\cite{liu2023security_region}提出了基于安全域的快速评估方法，但边界参数固定，无法自适应更新。Guo等\cite{guo2024tiered_limit}提出了运行方式聚类分档限额，但缺乏严格的安全域几何构造和闭环紧化机制。

针对上述不足，本文提出基于多输出高斯过程（MOGP）、贝叶斯优化（BO）与仿射内逼近（AIA）融合的闭环安全边界辨识框架，主要贡献如下：

\begin{enumerate}
    \item \textbf{建模层面}：在Kundur两区域系统中注入WECC标准REGCA1+REECA1+REPCA1新能源动态模型，构建8维参数空间与5级新能源渗透率场景。通过大规模仿真揭示新能源渗透率升高导致约束主导模式从"功角单一主导"向"多约束耦合"转变的物理机制。

    \item \textbf{方法层面}：提出"MOGP代理$\to$BO勘探$\to$AIA边界$\to$闭环紧化"的融合框架。多输出GP同时预测三约束严重度，BO以期望改进为采集函数定向搜索边界临界点，AIA通过凸包与LP分离超平面构造仿射安全域，闭环迭代逐步紧化边界并保证安全性。

    \item \textbf{验证层面}：在186个可行运行方式$\times$7种异构故障$\times$5级渗透率共6510次仿真中，全面验证了所提方法的代理模型精度（$R^2 > 0.85$）、勘探效率（评估次数减少40\%+）、边界安全性（安全率100\%）和工程效益（分档限额提升15\%--73\%）。
\end{enumerate}
